\section{Результаты и их обсуждение}
\subsection{Дисперсия золота и положение плазмонных полос}
Факторы деполяризации равны $L_\parallel=0{,}119339$ и
$L_\perp=0{,}440331$. Продольный максимум
$\operatorname{Im}\alpha$ прямого табличного расчёта находится при
1,7686 эВ, то есть на 0,2734 эВ ниже перехода КТ. Поперечный максимум
расположен при 2,3863 эВ. Таким образом, несущая 2,042 эВ оказывается
между двумя полосами, а не на продольном плазмонном резонансе.
Расстройка с продольным максимумом в три раза превышает спектральную
ширину интенсивности импульса. Уже это исключает перенос представления
о резонансном торцевом усилении без проверки комплексного локального поля.

\figone{fig01_material.png}{1}{Дисперсия поляризуемости золотого сфероида:
прямое табличное описание и представления с одним и 12 лоренцевыми
осцилляторами. Показаны действительная и мнимая части, обратная
поляризуемость и ошибки аппроксимации. Соответствует рисунку 1
вычислительного плана.}{fig:material}

В продольном направлении $N=12$ воспроизводит положение максимума
на используемой сетке; ширина полосы равна 0,07459 эВ против
0,07493 эВ в прямом расчёте. NRMS комплексной поляризуемости составляет
2,475\%, обратной поляризуемости --- 2,457\%. Для поперечного
направления эти ошибки равны 2,095 и 1,633\%, а максимум смещается
с 2,3863 до 2,3896 эВ. Следовательно, в пределах заданных критериев
получено согласованное многоосцилляторное представление амплитуды и фазы.

Есть существенное ограничение такого утверждения. Ширина широкой
поперечной полосы по глобальному основанию равна 1,2096 эВ для
табличного материала и 0,5106 эВ для $N=12$. Эта метрика чувствительна
к широкому фону, уровню полувысоты и межзонной дисперсии. Малый NRMS
не гарантирует малой ошибки любой нелокально определённой ширины.
Поэтому поперечную ширину из этой таблицы нельзя трактовать как
надёжно воспроизведённое время жизни отдельного плазмона. Для статьи
решающими служат дополнительные проверки непосредственно спектра КТ
вблизи 2,042 эВ.

Одноосцилляторная ветвь значительно грубее: NRMS поляризуемости равен
94,4\% для продольной и 36,1\% для поперечной ориентации.
Её продольный максимум на сетке находится при 0,989 эВ, а поперечный
выходит на верхнюю границу окна 3,5 эВ; корректная двусторонняя ширина
в этих случаях не выделяется. Эта ветвь сохраняется для количественной
оценки последствий \emph{конкретной} грубой подгонки во всём окне,
а не как оптимальная однорезонансная модель в узкой полосе около КТ.

\subsection{Усиление, спектральные сдвиги и уширение в пяти каналах}
Результаты FQS при $g=0{,}5$ нм сведены в табл.~\ref{tab:linear}.
Продольный боковой канал имеет наибольшее усиление и при фиксированной
энергии, и после оптимизации по энергии: $G_{\rm fixed}=3{,}116$,
$G_{\rm opt}=3{,}583$. Следующий по оптимизированному усилению канал
--- боковой радиальный, для которого эти величины равны 1,338 и 2,430.
Различие существенно: сдвиг радиальной линии уменьшает отклик именно
на несущей, хотя её собственный максимум остаётся высоким.

\begin{table}[!htbp]\centering\small
\caption{Линейные характеристики FQS при $g=0{,}5$ нм.
Сдвиг отсчитывается от максимума изолированной КТ;
$\Gamma_0=2{,}54142$ мэВ.}
\label{tab:linear}
\begin{tabular}{lrrrr}\toprule
Канал & $G_{\rm fixed}$ & $G_{\rm opt}$ & $\Delta E/\Gamma_0$ & $\Gamma/\Gamma_0$\\\midrule
Торцевой продольный & 0,648 & 0,883 & $-0,426$ & 1,426\\
Торцевой поперечный & 0,490 & 0,500 & $-0,073$ & 1,013\\
Боковой продольный & 3,116 & 3,583 & $-0,205$ & 1,057\\
Боковой радиальный & 1,338 & 2,430 & $-0,538$ & 1,192\\
Боковой тангенциальный & 0,143 & 0,156 & $-0,158$ & 1,024\\\bottomrule
\end{tabular}
\end{table}

Торцевая продольная конфигурация при этой энергии подавляет
монохроматический отклик, несмотря на сильную поляризуемость вдоль
длинной оси. Её линия заметно уширена, тогда как у наиболее эффективного
бокового продольного канала увеличение ширины составляет лишь 5,73\%.
Сильное возбуждающее поле и сильная обратная связь поэтому не являются
одним и тем же преимуществом: знаменатель $1-\beta K$ в
формуле~\eqref{eq:linear} может ослаблять выигрыш от числителя.
В FQS все пять линий при минимальном зазоре имеют красный сдвиг.
В частности, боковой продольный сдвиг равен примерно $-0{,}520$ мэВ,
а радиальный --- $-1{,}368$ мэВ. Для DD продольный боковой сдвиг,
напротив, положителен: $+0{,}212\Gamma_0$ вместо
$-0{,}205\Gamma_0$. Различие затрагивает даже знак эффекта.

\begin{figure}[!htbp]\centering
\includegraphics[width=.85\linewidth]{fig02a.png}\par
\includegraphics[width=.85\linewidth]{fig02b.png}
\caption{Зависимости оптимизированного усиления и сдвига экситонной
линии от зазора: FQS и DD, все пять каналов (панели 2a и 2b плана).
Оптимизация усиления проводится по одному общему энергетическому окну.}
\label{fig:gainshift}
\end{figure}
\begin{figure}[!htbp]\centering
\includegraphics[width=.85\linewidth]{fig02c.png}\par
\includegraphics[width=.85\linewidth]{fig02d.png}
\caption{Ширина линии относительно изолированной КТ и интегральное
расхождение спектров DD/FQS (панели 2c и 2d плана).
Штриховая линия на нижней панели соответствует $\delta_{\rm spec}=0{,}1$.
Треугольники у правой границы для продольных каналов отмечают только
последнюю принятую точку, а не доказанную границу непрерывной области.}
\label{fig:widtherror}
\end{figure}

Зависимость усиления от зазора не всегда монотонна.
Для бокового продольного канала максимум $G_{\rm fixed}$ на сетке
равен 3,2251 при $g=1$ нм; при 0,75 нм он практически такой же
(3,2245), а при 0,5 нм немного ниже. Для радиального канала максимум
фиксированного усиления 2,0406 достигается при 2 нм.
Однако $G_{\rm opt}$ обоих каналов наибольший при 0,5 нм.
Тем самым <<лучший зазор>> зависит от того, допускается ли перенастройка
энергии возбуждения.

Торцевой продольный $G_{\rm fixed}$ уменьшается до 0,0984 при 3 нм,
после чего возрастает до 0,6957 при 12 нм. Такой минимум согласуется
с интерференционным ослаблением суммарного поля; он не описывается
одним монотонным законом $R^{-6}$. Увеличение зазора в остальных
каналах также ведёт к ослаблению взаимодействия, но подход к
изолированной КТ может происходить как сверху, так и снизу.

\figone{fixed_and_pulse_gain.pdf}{1}{Дополнительное сопоставление FQS:
слева усиление на фиксированной энергии 2,042 эВ, справа отношение
слабых послеимпульсных заселённостей гибрида и изолированной КТ.
Логарифмическая вертикальная шкала показывает как усиление,
так и подавление. Это разные наблюдаемые; горизонтальная линия равна единице.}{fig:fixedpulse}

\subsection{Где дипольное приближение перестаёт быть количественным}
При $g=0{,}5$ нм спектральное расхождение $\delta_{\rm spec}$ равно
0,886; 0,637; 3,929; 2,084 и 0,942 в порядке каналов
табл.~\ref{tab:geometry}. Значения больше единицы допустимы:
метрика сравнивает ненормированные спектры, а DD сильно завышает
боковой отклик. В частности, $G_{\rm opt}$ бокового продольного
канала в DD равно 16,419 вместо 3,583, а радиального ---
7,754 вместо 2,430. Для торцевого продольного канала DD, напротив,
даёт 0,142 вместо 0,883. Направление ошибки зависит от геометрии.

С увеличением зазора в целом улучшается согласие, однако для торцевого
продольного канала $\delta_{\rm spec}$ достигает 2,232 при 3 нм:
именно здесь спектр FQS сильно подавлен и относительная ошибка
особенно чувствительна к компенсации полей. Нельзя устанавливать
границу DD по первому пересечению уровня 10\%, не проверив все
последующие узлы сетки.

\begin{table}[!htbp]\centering\small
\caption{Начало принятого хвоста сетки при допуске 10\%.
Для обоих критериев, $\delta_{\rm spec}$ и $\delta_F$, получены
одинаковые узлы, но критерии вычислялись независимо.
Один узел означает предварительную оценку на краю диапазона.}
\label{tab:DD}
\begin{tabular}{lrrl}\toprule
Канал & $g$, нм & Узлов в хвосте & Статус\\\midrule
Торцевой продольный & 12 & 1 & Только крайняя точка\\
Торцевой поперечный & 5 & 4 & Подтверждено на сетке\\
Боковой продольный & 12 & 1 & Только крайняя точка\\
Боковой радиальный & 10 & 2 & Подтверждено на сетке\\
Боковой тангенциальный & 7 & 3 & Подтверждено на сетке\\\bottomrule
\end{tabular}
\end{table}
При 12 нм $\delta_{\rm spec}$ продольных каналов всё ещё составляет
около 9,1\%. Поэтому утверждение об универсальной применимости DD
при зазоре, например, 5 нм не подтверждается. Таблица~\ref{tab:DD}
также не задаёт непрерывную границу между узлами и не распространяется
автоматически на другие частоты или формы частиц.

\subsection{Порог импульсного возбуждения и его связь со слабым откликом}
Порог изолированной КТ равен $61{,}672\ \uJ$, что соответствует
пиковой интенсивности $2{,}897$ ГВт/см$^2$ при принятой длительности.
Минимум FQS на рассчитанной сетке получается для бокового
продольного канала при $g=0{,}5$ нм:
\begin{equation}
\Fth^{\rm FQS}=16{,}075\ \uJ,\qquad
\frac{\Fth^0}{\Fth^{\rm FQS}}=3{,}8365.
\end{equation}
Соответствующая пиковая интенсивность составляет около
0,755 ГВт/см$^2$. Следующий канал, боковой радиальный, требует
$20{,}229\ \uJ$. Порог у торца при продольной поляризации
равен $49{,}229\ \uJ$, то есть выше бокового в 3,062 раза.

\begin{table}[!htbp]\centering\small
\caption{Пороговые флюенсы при $g=0{,}5$ нм, $\uJ$.
Обозначение $>200$ отражает статус ограничения диапазоном поиска;
оно не является вычисленным бесконечным порогом.}
\label{tab:threshold}
\begin{tabular}{lrrr}\toprule
Канал & FQS & DD & $\Fth^0/\Fth^{\rm FQS}$\\\midrule
Торцевой продольный & 49,229 & $>200$ & 1,253\\
Торцевой поперечный & 121,184 & 75,459 & 0,509\\
Боковой продольный & 16,075 & 3,168 & 3,837\\
Боковой радиальный & 20,229 & 7,574 & 3,049\\
Боковой тангенциальный & $>200$ & $>200$ & ---\\\bottomrule
\end{tabular}
\end{table}

\begin{figure}[!htbp]\centering
\includegraphics[width=.87\linewidth]{fig03a.png}\par
\includegraphics[width=.87\linewidth]{fig03b.png}
\caption{Нелинейный порог $\rho_{ee}(t_{\rm read})=0{,}5$,
нормированный на порог изолированной КТ, и его относительное
расхождение между DD и FQS (панели 3a, 3b плана).
Неразрешённые пороги не превращаются в конечные отношения.
Показанные линии соединяют рассчитанные точки.}
\label{fig:threshold}
\end{figure}

Для бокового продольного канала порог FQS монотонно возрастает
от 16,075 при 0,5 нм до 17,579 при 1 нм, 24,262 при 3 нм
и 47,397 $\uJ$ при 12 нм. При переходе от 0,5 к 1 нм
требуемый флюенс возрастает лишь на 9,35\%.
Следовательно, точка 0,5 нм является минимумом \emph{на сетке};
она не доказывает существование физически оптимального контакта
при ещё меньшем расстоянии. Для радиального канала порог растёт
от 20,229 до 52,532 $\uJ$ на том же диапазоне.

Торцевой продольный канал демонстрирует немонотонное поведение.
Порог возрастает до 174,990 $\uJ$ при 1,5 нм; при 2, 3 и 5 нм
он не разрешён в исследованном диапазоне. При 7 нм снова получается
конечное значение 148,976, при 12 нм --- 88,620 $\uJ$.
Эта последовательность согласуется с минимумом линейного отклика
в промежуточной области. Боковой тангенциальный канал имеет
неразрешённый порог при 0,5--1 нм, после чего его порог убывает
от 190,460 при 1,5 нм до 67,776 $\uJ$ при 12 нм.

Отдельно следует объяснить, почему торцевой продольный
$G_{\rm fixed}<1$, но его импульсный порог при 0,5 нм ниже
изолированного. Импульс охватывает конечную полосу частот и возбуждает
последующую связанную динамику; его действие не сводится к значению
$S$ в одной точке. Отношение слабых послеимпульсных заселённостей
для этого канала равно 1,250, тогда как $G_{\rm fixed}=0,648$.
Для продольного бокового канала слабый импульсный выигрыш равен
3,784 и близок, но не тождествен пороговому выигрышу 3,837.
Для радиального канала соответствующие значения составляют 2,875 и 3,049.

Наименьший порог DD равен $3{,}168\ \uJ$; для той же
конфигурации FQS требует в 5,075 раза больший флюенс.
Относительное занижение порога равно 80,29\%.
Радиальная ошибка составляет 62,56\%, торцевая поперечная ---
37,73\%. При больших зазорах расхождение порогов уменьшается,
но отсутствие значения в одной модели исключает вычисление
$\delta_F$ для соответствующей пары.

\subsection{Почему боковой продольный канал эффективнее торцевого}
Проверка геометрического механизма проведена независимо от поиска
нелинейных порогов. При 2,042 эВ таблица золота даёт
$\varepsilon_{\rm Au}=-9{,}9506+1{,}4634i$.
Для голой частицы и $g=0{,}5$ нм точное электростатическое поле
в центре будущей КТ определяется коэффициентами
\begin{equation}
B_{\rm tip}\simeq-2{,}1103+0{,}7367i,\qquad
B_{\rm side}\simeq0{,}9703-0{,}3387i.
\end{equation}
Хотя $|B_{\rm tip}|^2=4{,}996$ больше
$|B_{\rm side}|^2=1{,}056$, полное возбуждающее поле содержит
падающую волну:
\begin{equation}
|1+B|^2=1+2\operatorname{Re}B+|B|^2.
\end{equation}
В результате $|1+B_{\rm tip}|^2=1{,}7755$, а
$|1+B_{\rm side}|^2=3{,}9970$. У торца интерференционный член
отрицателен, сбоку положителен. Сильнее рассеянное поле не означает
сильнее суммарное поле.

Обратная связь дополнительно различает эти каналы. При той же энергии
$1-\beta K\simeq1{,}4355-0{,}8236i$ у торца и
$1{,}0573-0{,}4061i$ сбоку. Деление $|1+B|^2$ на квадрат
модуля соответствующего знаменателя воспроизводит
$G_{\rm fixed}=0{,}6483$ и 3,1157.
Таким образом, наблюдаемая иерархия прослеживается от фазы внешнего поля
до самосогласованного отклика КТ, а не вводится как эмпирическое
правило выбора положения.

\figone{tip_side_fields.png}{1}{Независимая проверка локального поля
продольной яркой моды у торца и сбоку. Расстояние на профилях
отсчитывается от поверхности до точки наблюдения. При номинальном
зазоре центр КТ находится на расстоянии 2,5 нм от поверхности.
Поле голого металла рассматривается отдельно от обратного действия КТ.}{fig:fields}

Непосредственно на поверхности продольное усиление интенсивности
у торца действительно достигает 153,05 против 7,66 сбоку.
Однако профили пересекаются уже на расстоянии 1,824 нм от поверхности,
а центр КТ при $g=0{,}5$ нм находится на расстоянии 2,5 нм.
Торцевой максимум на поверхности поэтому не является полем,
действующим на принятый точечный переход.
На продольном плазмонном максимуме 1,7686 эВ картина поля иная:
в тех же центрах усиления составляют около 483,6 и 102,1.
Эти числа описывают только голую частицу. Они не предсказывают
нелинейный порог КТ с переходом 2,042 эВ при переносе одной лишь несущей.

\subsection{Как аппроксимация материала меняет наблюдаемые}
Прямой табличный расчёт и $N=12$ практически совпадают в области
экситонной линии на рис.~\ref{fig:materialresponse}.
При $g=0{,}5$ нм NRMS спектра составляет 0,3233\% для
торцевого продольного, 0,0595\% для торцевого поперечного,
0,0534\% для бокового продольного, 0,2715\% для радиального
и 0,2840\% для тангенциального канала. На трёх контрольных
зазорах 0,5, 10 и 12 нм максимальная ошибка положения пика равна
$0{,}001702\Gamma_0$, ширины --- 0,0337\%, усиления --- 0,2724\%.
Эти ошибки значительно меньше заданных допусков конечных наблюдаемых.

\figone{fig04.png}{1}{Спектр КТ при прямом табличном материале,
аппроксимациях с одним и 12 осцилляторами для двух ведущих
каналов, $g=0{,}5$ нм (рисунок 4 плана). Нижние панели показывают
разности относительно прямого спектра с нормировкой на его максимум.
Принятие многоосцилляторной ветви проверено по всем пяти каналам.}{fig:materialresponse}

При $N=1$ соответствующие NRMS спектра составляют
34,9; 21,6; 72,4; 40,7 и 110,0\%. В частности,
оптимизированное усиление бокового продольного канала занижается
примерно на 73,4\%. Пассивная и устойчивая аппроксимация тем самым
может оказаться непригодной для количественного описания наблюдаемой.

\begin{table}[!htbp]\centering\small
\caption{Влияние материальной модели на порог FQS при $g=0{,}5$ нм.
Значения получены из отдельной плотной сетки зависимостей от флюенса
(рисунок 5 плана); малая разница с табл.~\ref{tab:threshold} обусловлена
процедурой извлечения порога. Флюенсы в $\uJ$.}
\label{tab:materialthreshold}
\begin{tabular}{lrrr}\toprule
Канал & $N=1$ & $N=12$ & Отношение $1/12$\\\midrule
Изолированная КТ & 61,673 & 61,673 & 1,000\\
Торцевой продольный & 38,912 & 49,234 & 0,790\\
Торцевой поперечный & 97,279 & 121,184 & 0,803\\
Боковой продольный & 53,498 & 16,075 & 3,328\\
Боковой радиальный & 28,896 & 20,228 & 1,429\\
Боковой тангенциальный & 172,624 & $>200$ & ---\\\bottomrule
\end{tabular}
\end{table}

\figone{fig05.png}{1}{Послеимпульсная заселённость и её разность
для $N=1$ и $N=12$ при увеличении флюенса, $g=0{,}5$ нм
(рисунок 5 плана). Показаны боковые продольный и радиальный каналы.
Кружками отмечены первые пересечения уровня 0,5; последующее
уменьшение заселённости относится к когерентной нелинейной динамике.}{fig:materialfluence}

В нелинейной задаче отличие ещё существеннее, чем просто изменение
масштаба кривой. При $N=1$ радиальный канал имеет меньший порог,
чем продольный боковой (28,90 против 53,50), тогда как при $N=12$
порядок обратный (20,23 против 16,08 $\uJ$).
Одноосцилляторное описание может изменить выбор геометрии.
На полной сетке флюенсов максимальная абсолютная разность
заселённостей продольного бокового канала достигает 0,908;
часть этой разности связана со смещением рабиевского максимума,
а не с постоянным коэффициентом усиления.
Ветка $N=1$ здесь подгоняется к одиночному сфероиду и не тождественна
эффективному плазмонному осциллятору димера из~\cite{shah2013}.

\subsection{Динамика при одинаковом падающем импульсе}
Временные зависимости сравниваются при общем флюенсе
$16{,}07548\ \uJ$, выбранном по порогу $N=12$
продольного бокового канала. Меняется геометрия или приближение,
но амплитуда падающего импульса не перенормируется отдельно для
каждой кривой. Результаты приведены в табл.~\ref{tab:dynamics}
и на рис.~\ref{fig:dynamicsnear}, \ref{fig:dynamicsfar}.

\begin{table}[!htbp]\centering\small
\caption{Заселённость при $t_{\rm read}=3{,}6$ пс и общем флюенсе
$16{,}07548\ \uJ$.}
\label{tab:dynamics}
\begin{tabular}{lrr}\toprule
Система & $g=0{,}5$ нм & $g=12$ нм\\\midrule
Изолированная КТ & 0,152966 & 0,152966\\
FQS, боковой продольный & 0,500005 & 0,195771\\
FQS, боковой радиальный & 0,411485 & 0,177826\\
DD, боковой продольный & 0,793060 & 0,212234\\
DD, боковой радиальный & 0,828376 & 0,183640\\\bottomrule
\end{tabular}
\end{table}

При малом зазоре DD сильно завышает конечное возбуждение и даёт
обратный порядок двух заселённостей при данном фиксированном флюенсе.
Это не противоречит меньшему продольному порогу DD: одна и та же
амплитуда помещает разные каналы на разные участки рабиевской кривой.
Для продольного DD максимальная заселённость достигает 0,9616
уже при $t=10{,}12$ фс, затем уменьшается до 0,7931.
Для продольного FQS максимум 0,5127 наблюдается около 41,95 фс
и к моменту считывания уменьшается до 0,5000. Радиальный FQS
достигает максимума 0,4525 около 38,01 фс и заканчивает на 0,4115.
Временное смещение и изменение формы отклика показывают, что
замена пространственного ядра меняет динамику, а не только
стационарную амплитуду.

\figone{fig06_near_pulse_zoom.png}{1}{Динамика около импульса при
$g=0{,}5$ нм: заселённость и нормированное падающее поле
(увеличенный фрагмент рисунка 6 плана).
Две геометрии FQS сравниваются с DD и изолированной КТ.}{fig:dynamicsnear}
\figone{fig06_far_pulse_zoom.png}{1}{То же сравнение при $g=12$ нм.
Флюенс оставлен равным порогу ближнего продольного бокового канала,
поэтому ни одна из дальних гибридных кривых не обязана достигать 0,5.}{fig:dynamicsfar}

При 12 нм DD и FQS сближаются: относительная разность конечных
заселённостей равна 8,41\% для продольного и 3,27\% для радиального
канала. Тем не менее металл ещё влияет на результат: продольная
заселённость FQS превышает изолированную примерно в 1,28 раза.
Сходимость DD и FQS поэтому не следует путать с полным исчезновением
взаимодействия. Полные временные интервалы показаны в приложении.

\subsection{Нелинейный спектр работы поля: результат S1}
Спектры рассчитывались для бокового продольного канала при
$g=0{,}5$ нм и трёх общих флюенсах:
$0{,}0005$, $16{,}07548$ и $59{,}97128\ \uJ$.
Первый соответствует слабому полю, второй --- порогу $N=12$,
третий --- выбранному узлу около первого нелинейного максимума
этой ветви. Последний не объявляется точно оптимизированной амплитудой.
Окно наблюдения после центра импульса составляет 12 пс,
спектральная сетка --- 8001 точку в интервале 1,90--2,18 эВ.

\figone{supp01.png}{1}{Спектры $\sigma_{\rm QS,work}$ и разности
$\Delta\sigma$ относительно голого металла для трёх флюенсов
и двух материальных представлений (S1 плана). Флюенсы в обеих
ветвях одинаковы, поэтому пороговой амплитуде $N=12$ не соответствует
порог $N=1$. Отрицательная разность с фоном не означает отрицательной
полной работы поля.}{fig:work}

В окрестности $2{,}042\pm0{,}010$ эВ максимум добавки
$\Delta\sigma$ для $N=12$ уменьшается от
$3{,}0245\cdot10^{-13}$ до $1{,}9548\cdot10^{-13}$ и затем до
$1{,}5004\cdot10^{-14}$ см$^2$. В сильном выбранном импульсе
минимум достигает $-1{,}6353\cdot10^{-14}$ см$^2$:
экситонная добавка существенно ослабевает и приобретает участок
с подавлением относительно голого металла. Для $N=1$ положительные
максимумы в том же окне равны $7{,}4436$, $6{,}7751$ и
$4{,}8923\cdot10^{-14}$ см$^2$. Следовательно, грубое описание
материала меняет как исходный контраст, так и степень его
нелинейного подавления.

Полная работа поля остаётся положительной во всех шести расчётах.
Для $N=12$ она равна $4{,}3021\cdot10^{-22}$,
$1{,}3793\cdot10^{-17}$ и $5{,}1183\cdot10^{-17}$ Дж.
Её отношение к флюенсу меняется сравнительно слабо:
от $8{,}6041\cdot10^{-13}$ до $8{,}5346\cdot10^{-13}$ см$^2$,
при значении голой частицы $8{,}4889\cdot10^{-13}$ см$^2$.
Резкое изменение узкой спектральной структуры сочетается с малым
изменением интегрального отклика, в котором преобладает металл.
Это объясняет необходимость анализировать одновременно полный спектр,
вычитание фона и интегральную работу.

Наблюдаемое изменение асимметричной добавки согласуется с нелинейным
изменением когерентного отклика КТ. Однако сами эти кривые не
доказывают воспроизведение сверхрассеяния или конкретного механизма
обращения резонанса Фано из~\cite{shah2013}: отличаются геометрия,
расстройка, затухание заселённости и определение наблюдаемой.
Для такого вывода потребовались бы отдельный фазовый анализ и
согласованное радиационное описание.
